% =====================================================================
%  Seance 3 -- A. cours 1 : detection d'anomalies non supervisee (frames 1--9)
% =====================================================================
\section{Séance 3 --- Modèles d'apprentissage automatique appliqués aux infrastructures}
\chapterslide{Séance 3}{Modèles d'apprentissage automatique appliqués aux infrastructures}
{4 h --- 2 h cours / 2 h TP --- Partie 3 du syllabus EPD-GIT-5534}
{\textbf{Objectifs :} détecter des anomalies \textbf{sans étiquette}
(IsolationForest, LOF) et choisir un seuil $\bullet$ comprendre pourquoi
l'\textbf{exactitude est trompeuse} à 0,35 \% de positifs $\bullet$ entraîner
et comparer 4 modèles supervisés (régression logistique, arbre, forêt,
gradient boosting) sur \file{nosys\_features.csv} $\bullet$ lire une matrice
de confusion, une courbe précision-rappel, et choisir un seuil \textbf{métier}
$\bullet$ valider un modèle par validation croisée \textbf{temporelle}.}

\begin{frame}{Plan de la séance (4 h)}
\begin{columns}[t]
\begin{column}{0.48\textwidth}
\begin{kbox}[Cours magistral --- 2 h]
\begin{enumerate}
\item \textbf{Anomalies sans étiquette} (40 min) : z-score/IQR (rappel),
IsolationForest, LOF --- scores et seuils.
\item \textbf{Classification sur classe rare} (45 min) : 4 familles de
modèles, précision/rappel/F1, courbe précision-rappel, coût asymétrique.
\item \textbf{Choisir le modèle} (35 min) : validation croisée
temporelle, compromis interprétabilité/performance.
\end{enumerate}
\end{kbox}
\end{column}
\begin{column}{0.48\textwidth}
\begin{ebox}[Travaux pratiques --- 2 h]
\textbf{« Détecter les incidents avant la panne »} sur
\file{nosys\_features.csv} :
\begin{enumerate}
\item B1 --- IsolationForest sans étiquette (25 min)
\item B2 --- recouvrement détections / alarmes réelles (20 min)
\item B3 --- forêt balanced + matrice de confusion (30 min)
\item B4 --- courbe PR et seuil métier $\le$ 5 alertes/jour/serveur (30 min)
\item B5 --- livrable : tableau comparatif + recommandation (15 min)
\end{enumerate}
\end{ebox}
\end{column}
\end{columns}
\vspace{1mm}
\begin{center}\small
Matière première : \textbf{votre} \file{nosys\_features.csv} du TP2 ---
34 560 fenêtres, dont \textbf{120 incidents} (0,35 \%).
\end{center}
\end{frame}

\begin{frame}{Rappel séance 2 --- et ses limites}
\begin{columns}[t]
\begin{column}{0.5\textwidth}
\begin{block}{Les seuils statistiques de la séance 2}
\begin{itemize}
\item \textbf{z-score} : flag si $|x - \mu| > 3\sigma$ (par variable).
\item \textbf{IQR} : flag si hors $[Q_1 - 1{,}5 \times \text{IQR}\,;$
$Q_3 + 1{,}5 \times \text{IQR}]$.
\item Appliqués ici à \py{lat\_p99} sur le vrai dataset :
\textbf{72} dépassements z-score, \textbf{82} dépassements IQR
(chiffres mesurés, f16).
\end{itemize}
\end{block}
\begin{wbox}[Trois limites fatales en production]
\begin{enumerate}
\item \textbf{Une variable à la fois} : un incident = une
\emph{combinaison} (CPU normal + latence haute + perte).
\item \textbf{Seuil fixé à la main} : 3$\sigma$, pourquoi pas 2,5 ?
\item \textbf{Aucune notion de rareté jointe} : le 99\% percentile
devient une alerte permanente.
\end{enumerate}
\end{wbox}
\end{column}
\begin{column}{0.46\textwidth}
\begin{center}
\includegraphics[width=\textwidth]{f16_seuils_statistiques.png}
\end{center}
\end{column}
\end{columns}
\end{frame}

\begin{frame}{IsolationForest : isoler plutôt que modéliser}
\begin{columns}[t]
\begin{column}{0.5\textwidth}
\begin{block}{Le principe en une phrase}
Une observation \textbf{rare} est isolée par \textbf{peu de coupes
aléatoires} ; une observation normale demande beaucoup de coupes.
\begin{itemize}
\item 300 arbres coupent l'espace au hasard ; la \textbf{profondeur
moyenne} d'isolement donne le \textbf{score d'anormalité}.
\item Aucune étiquette requise : apprentissage \textbf{non supervisé}.
\item \py{IsolationForest(contamination=0.01)} : contamination $=$
fraction d'anomalies \emph{attendue} --- elle fixe le seuil de décision.
\end{itemize}
\end{block}
\begin{kbox}[Pourquoi ça marche pour l'infra]
Le volume $(cpu, ram, latence, perte, temp)$ d'une fenêtre pré-incident
est \textbf{loin de tout le nuage} de fonctionnement normal : isolable
en peu de coupes, sans jamais avoir vu d'incident étiqueté.
\end{kbox}
\end{column}
\begin{column}{0.46\textwidth}
\begin{center}
\includegraphics[width=\textwidth]{f17_isolation_forest.png}
\end{center}
\end{column}
\end{columns}
\end{frame}

\begin{frame}[fragile]{IsolationForest en 8 lignes}
\begin{codebox}[Le squelette à retenir]
\begin{lstlisting}
from sklearn.ensemble import IsolationForest

iso = IsolationForest(n_estimators=300,   # arbres
                      contamination=0.01, # seuil implicite
                      random_state=5534, n_jobs=2)
iso.fit(X_train_sans_etiquette)

score = -iso.score_samples(X)   # grand = anormal
anomalies = X[score > score[train].quantile(0.99)]
\end{lstlisting}
\end{codebox}
\begin{columns}[t]
\begin{column}{0.48\textwidth}
\begin{block}{Les 3 réglages qui comptent}
\begin{itemize}
\item \py{n\_estimators} : 300 suffit, au-delà score stable.
\item \py{contamination} : votre hypothèse de rareté --- à
justifier par le métier (ici $\approx$ 1 \%).
\item \py{random\_state=5534} : reproductibilité, comme partout.
\end{itemize}
\end{block}
\end{column}
\begin{column}{0.48\textwidth}
\begin{wbox}[Deux erreurs de débutant]
\begin{enumerate}
\item Entraîner sur \textbf{train+test} : le score du test est
alors contaminé par sa propre distribution.
\item Oublier \py{fillna} : \py{cpu\_delta} de la première fenêtre
est NaN --- sklearn refuse les NaN.
\end{enumerate}
\end{wbox}
\end{column}
\end{columns}
\end{frame}

\begin{frame}{Ce que l'IsolationForest trouve \emph{vraiment} (mesuré)}
\begin{columns}[t]
\begin{column}{0.52\textwidth}
\begin{ebox}[Résultats réels sur nosys\_features.csv]
Entraîné sur le train (24 192 fenêtres, \textbf{sans} la colonne
\py{incident\_30min}) :
\begin{itemize}
\item les \textbf{20 fenêtres les plus anormales} sont
\textbf{20/20} des pré-incidents (f17a) ;
\item mais au seuil « 1 \% du train », le test contient
\textbf{174 anomalies} dont seulement \textbf{40 incidents}
--- \textbf{134 fausses alarmes} (f17b).
\end{itemize}
\end{ebox}
\begin{wbox}[Piège n°1 : anomalie $\ne$ incident]
Une \textbf{anomalie} est un point statistiquement rare. Un
\textbf{incident} est une anomalie \emph{qui casse le service}.
Backup de nuit, pic de batch, déploiement : anormaux, bénins.
L'IsolationForest \textbf{name} le soupçon, il ne \textbf{classe} pas
la panne. C'est le travail du supervisé (cours 2).
\end{wbox}
\end{column}
\begin{column}{0.44\textwidth}
\begin{block}{LOF : la densité locale (complément)}
Le LOF (Local Outlier Factor) compare la densité \textbf{autour} d'un
point à celle de ses voisins : utile quand les clusters normaux ont des
\textbf{densités très différentes} (nuit vs pic de journée).
\begin{itemize}
\item LOF $\approx 1$ : densité normale ; LOF $\gg 1$ : hors densité.
\item Sur nosys, IsolationForest et LOF pointent essentiellement les
mêmes fenêtres --- on retient IsolationForest, scalable et parallèle
(\py{n\_jobs=2}).
\end{itemize}
\end{block}
\end{column}
\end{columns}
\end{frame}

\begin{frame}{Scores, seuils : qui décide de « alerte » ?}
\begin{columns}[t]
\begin{column}{0.48\textwidth}
\begin{block}{La chaîne complète}
score brut $s(x)$ $\Rightarrow$ seuil $t$ $\Rightarrow$ décision
$s(x) \ge t$.
\begin{itemize}
\item \textbf{z-score / IQR} : seuil implicite (3$\sigma$, 1,5 IQR).
\item \textbf{IsolationForest} : seuil $=$ quantile du score sur le
train (1 \% ici) ou paramètre \py{contamination}.
\item \textbf{Supervisé} : seuil sur la probabilité prédite ---
cours 2 et TP B4.
\end{itemize}
\end{block}
\end{column}
\begin{column}{0.48\textwidth}
\begin{kbox}[Le seuil est un choix MÉTIER, pas statistique]
Descendre le seuil $=$ plus de rappel, plus de fatigue NOC.
Le monter $=$ moins de bruit, incidents manqués.
Il n'existe \textbf{pas} de seuil optimal universel : il y a un seuil
par \textbf{budget d'opérateurs}. Le TP B4 le chiffre : 5 alertes
/jour/serveur maximum.
\end{kbox}
\end{column}
\end{columns}
\end{frame}

\begin{frame}{Quiz 1 --- anomalies non supervisées (5 min)}
\begin{qbox}[3 questions]
\begin{enumerate}
\item Une fenêtre anormale selon IsolationForest est :\\
(a) forcément un incident \quad (b) statistiquement rare --- à vérifier
contre les alarmes réelles \quad (c) une panne confirmée
\item Le paramètre \py{contamination} sert à :\\
(a) initialiser les arbres \quad (b) fixer la fraction d'anomalies
attendue, donc le seuil \quad (c) paralléliser
\item Pourquoi le z-score suffit-il rarement seul ?\\
(a) trop lent \quad (b) une variable à la fois, seuil arbitraire,
ignore les combinaisons \quad (c) il exige des étiquettes
\end{enumerate}
\end{qbox}
\vspace{1mm}
\begin{solbox}[Réponses]
\textbf{1 : b} \quad \textbf{2 : b} \quad \textbf{3 : b} --- d'où le
passage au modèle joint et supervisé.
\end{solbox}
\end{frame}
